%----------------------------------------------------------------------------------------
%	PACKAGES AND OTHER DOCUMENT CONFIGURATIONS
%----------------------------------------------------------------------------------------

\documentclass[a4paper,12pt]{memoir} % Font and paper size

\input{struc.tex} % Include the file specifying document layout and packages

%----------------------------------------------------------------------------------------
%	NAME AND CONTACT INFORMATION 
%----------------------------------------------------------------------------------------

\userinformation{
\begin{flushright}
\includegraphics[width=0.7\columnwidth,height=80pt]{photo.jpg}\\[\baselineskip] % Your photo

\small
Philipp Hanussek \\
\url{hanup@mailbox.org} \\
575 854 652 \\
\Sep
\textbf{Adres} \\
ul. Bytomska 105 \\
42-672 Wieszowa \\
\Sep
\textbf{Języki} \\
Niemiecki (ojczysty) \\
Angielski (C2, IELTS - styczeń 2024) \\
Francuski (C2) \\
Polski (C1) \\
\vfill
\end{flushright}
}
%----------------------------------------------------------------------------------------

\begin{document}

\userinformation % Print your information in the left column

\framebreak % End of the first column

%----------------------------------------------------------------------------------------
%	HEADING
%----------------------------------------------------------------------------------------
%----------------------------------------------------------------------------------------
\cvheading{Philipp Hanussek}

\cvsubheading{Student Informatyki / Badania ML}

\CVSection{Umiejętności techniczne}

\CVItem{Praca}{
Python do uczenia maszynowego (PyTorch, NumPy/CuPy, Pandas);
frameworki ML (Optuna, Streamlit); ABAP; REST API; SQL; Bash
}

\CVItem{Projekty}{
Scikit-learn; Julia; C do obliczeń wysokowydajnych (wielowątkowość);
Matlab; dokumentacja techniczna (UML); ITIL4
}

\Sep

%----------------------------------------------------------------------------------------
% EDUCATION
%----------------------------------------------------------------------------------------

\CVSection{Wykształcenie}

\CVItem{09/2024 -- 09/2026, Sorbonne Université (Paryż, Francja)}{
Magister Informatyki \\
Kursy: Deep Learning, High Performance Computing, kryptografia (w tym kwantowa), sieci neuronowe do rozpoznawania obrazów
}

\CVItem{2020 -- 2024, FH Aachen (Akwizgran, Niemcy)}{
Licencjat z Informatyki \\
Ocena końcowa: GPA 3.5/4 (z wyróżnieniem) \\
Praca dyplomowa: Comparison of Factoring Methods on the D-Wave Quantum Annealer \\
Osiągnięcia:
\begin{itemize}
\item Top 5\% rocznika (2021, 2022)
\item Stypendium za wyniki w nauce (2023)
\end{itemize}
}

\Sep

%----------------------------------------------------------------------------------------
% EXPERIENCE
%----------------------------------------------------------------------------------------

\CVSection{Doświadczenie}

\CVItem{03/2026 -- 08/2026, \textit{Badacz kwantowy}, Polska Akademia Nauk (Gliwice, Polska)}{
\begin{itemize}
\item Ocena wykorzystania współczesnych urządzeń kwantowych i klasycznych (GPU) do trenowania generatywnych sieci neuronowych
\item Główny autor publikacji naukowej przyjętej przez American Physical Society dotyczącej symulacji układów kwantowych
\end{itemize}
}

\CVItem{03/2024 -- 07/2024, \textit{Asystent studencki}, Jülich Supercomputing Centre, Forschungszentrum Jülich (Niemcy)}{
\begin{itemize}
\item Udział w projekcie oceniającym możliwości współczesnych komputerów kwantowych w rozwiązywaniu problemów kryptograficznych 
\item Opracowanie pipeline’ów symulacyjnych i eksperymentalnych do benchmarkowania systemów kwantowych
\end{itemize}
}

\CVItem{04/2022 -- 05/2024, \textit{Programista automatyzacji (część etatu)}, Zurich Insurance Group (Kolonia, Niemcy)}{
\begin{itemize}
\item Projektowanie i wdrażanie rozwiązań automatyzujących procesy płatności (ABAP, skrypty)
\item Odpowiedzialność za migrację rozwiązań automatyzacji do SAP FS-CD
\item Opracowanie prototypowego dashboardu SAP UI5/Fiori i integracja z backendem w ramach globalnego zespołu IT w Barcelonie
\end{itemize}
}

\CVItem{09/2021 -- 07/2022, \textit{Tutor studencki}, FH Aachen}{
\begin{itemize}
\item Prowadzenie zajęć przygotowawczych dla ponad 20 studentów
\item Wsparcie w zaawansowanych zagadnieniach matematycznych
\end{itemize}
}

\Sep

%------------------------------------------------

\Sep % Extra whitespace after the end of a major section


%------------------------------------------------



%----------------------------------------------------------------------------------------
%	NEW PAGE DELIMITER
%----------------------------------------------------------------------------------------

\clearpage % Start a new page

\begin{flushright}

\end{flushright}

\framebreak % End of the first column

%----------------------------------------------------------------------------------------
% PUBLICATIONS
%----------------------------------------------------------------------------------------

\CVSection{Publikacje}

\CVItem{2026}{
\href{https://journals.aps.org/prapplied/abstract/10.1103/rrf3-jm5m}
{\textit{Quantum-inspired dynamical models on quantum and classical annealers}} \\
P. Hanussek, J. Pawłowski, Z. Mzaouali, B. Gardas \\
Phys. Rev. Applied
}

\CVItem{2024}{
\href{https://arxiv.org/abs/2410.14397}
{\textit{The State of Factoring on Quantum Computers}} \\
D. Willsch, P. Hanussek, G. Hoever, M. Willsch, F. Jin, H. De Raedt, K. Michielsen \\
NIC Symposium 2025
}
\Sep

%----------------------------------------------------------------------------------------
% VOLUNTEERING
%----------------------------------------------------------------------------------------

\CVSection{Wolontariat}

\CVItem{09/2025 -- obecnie, Paryż / Gliwice}{
Wolontariusz \\
Cotygodniowa działalność (2h): dystrybucja posiłków dla osób w potrzebie, wsparcie i kierowanie do lokalnych usług pomocowych
}

\Sep
\end{document}